\subsection{Valenzelektronen und chemische Bindung}
\label{subsec:valenzelektronen-chemische-bindung}

Im vorherigen Abschnitt wurde deutlich, dass die Struktur der Elektronenhülle
die Ordnung des Periodensystems bestimmt. Für die chemischen Eigenschaften
eines Atoms sind dabei nicht alle Elektronen gleich wichtig. Entscheidend sind
vor allem die Elektronen der äußersten besetzten Schale. Sie werden
Valenzelektronen\index{Valenzelektronen} genannt.

Valenzelektronen bestimmen, wie ein Atom mit anderen Atomen wechselwirkt. Sie
entscheiden, ob ein Atom Elektronen leicht abgibt, Elektronen aufnimmt oder
Elektronen mit anderen Atomen teilt. Damit bilden sie die Grundlage für die
chemische Bindung\index{chemische Bindung}.

Ein Atom mit nur wenigen Valenzelektronen kann diese oft relativ leicht
abgeben. Dies ist typisch für viele Metalle. Ein Atom, dem nur wenige
Elektronen zur vollständig besetzten äußeren Schale fehlen, nimmt dagegen
leicht Elektronen auf oder teilt Elektronen mit anderen Atomen. Dies ist
typisch für viele Nichtmetalle.

\begin{PhysicsBox}[Valenzelektronen]
	\label{box:valenzelektronen}
	Valenzelektronen sind die Elektronen der äußersten besetzten Schale eines
	Atoms. Sie bestimmen maßgeblich die chemischen Eigenschaften eines Elements.
	
	Sie können
	\begin{itemize}
		\item abgegeben werden,
		\item aufgenommen werden,
		\item oder mit anderen Atomen geteilt werden.
	\end{itemize}
	
	Aus diesen Möglichkeiten entstehen die wichtigsten Arten chemischer Bindung.
\end{PhysicsBox}

Die chemische Bindung lässt sich anschaulich als Versuch verstehen, eine
energetisch günstigere Elektronenanordnung zu erreichen. Besonders stabil sind
Elektronenkonfigurationen, bei denen die äußere Schale vollständig besetzt ist.
Diese Situation findet man bei den Edelgasen\index{Edelgase}. Viele Atome
gehen Bindungen ein, weil sie dadurch eine ähnliche äußere Elektronenstruktur
wie ein Edelgas erreichen können.

Diese Vorstellung wird häufig als Edelgasregel\index{Edelgasregel} oder
Oktettregel\index{Oktettregel} bezeichnet. Sie ist eine nützliche Näherung,
besonders für viele Hauptgruppenelemente. Dabei geht es darum, dass Atome in
vielen Verbindungen acht Elektronen in ihrer äußeren Schale erreichen. Für
Wasserstoff und Helium gilt eine kleinere stabile Konfiguration mit zwei
Elektronen.

Die Oktettregel darf jedoch nicht als unveränderliches Naturgesetz verstanden
werden. Sie ist ein einfaches Modell, das viele chemische Bindungen gut
erklärt, aber nicht alle. Besonders bei Übergangsmetallen, größeren Atomen und
komplexeren Molekülen reichen einfache Schalenbilder nicht mehr aus. Dort
muss man die Elektronen mit Hilfe von Orbitalen und quantenmechanischen
Zuständen genauer beschreiben.

\begin{DidacticBox}[Warum Atome Bindungen eingehen]
	\label{box:warum-atome-bindungen-eingehen}
	Atome gehen Bindungen ein, wenn dadurch eine energetisch günstigere
	Elektronenanordnung entsteht.
	
	Anschaulich kann man sagen: Viele Atome erreichen durch Abgabe, Aufnahme oder
	gemeinsames Nutzen von Valenzelektronen eine stabilere äußere
	Elektronenkonfiguration.
\end{DidacticBox}

Damit wird klar, warum das Elektron in der Chemie eine zentrale Rolle spielt.
Chemische Bindungen sind keine rein mechanischen Verbindungen zwischen Atomen.
Sie entstehen aus der Wechselwirkung ihrer Elektronenhüllen. Die äußeren
Elektronen bestimmen, ob ein Atom reaktionsträge ist, leicht Ionen bildet oder
stabile Moleküle eingeht.

Im nächsten Abschnitt werden die wichtigsten Bindungsarten genauer betrachtet:
Ionenbindung, kovalente Bindung und Metallbindung.